\documentclass[11pt,a4paper]{article}
% Shared preamble for RuPhO English problem statements (match T1 2026 style).
% Use: \input{latex/rupho_en_preamble.tex} after \documentclass.
\usepackage[margin=1in]{geometry}
\usepackage{amsmath,amssymb}
\usepackage{graphicx}
\usepackage{float}
\usepackage{enumitem}
\usepackage{microtype}
\usepackage{caption}
\usepackage{tabularx}
\usepackage{hyperref}

\captionsetup{font=small,labelfont=bf,skip=6pt}

\setlist[itemize]{leftmargin=1.2cm}
\setlist[enumerate]{leftmargin=1.2cm}

% One outer box per subproblem: [id | statement | points] (no inner rules)
\newcommand{\problembox}[3]{%
  \par\addvspace{0.42em}%
  \noindent
  \setlength{\fboxsep}{7.5pt}%
  \setlength{\fboxrule}{0.95pt}%
  \fbox{%
    \begin{minipage}{\dimexpr\linewidth-2\fboxsep-2\fboxrule\relax}
    \begin{tabularx}{\linewidth}{@{}>{\raggedright\arraybackslash\bfseries}p{2.1em} @{\hspace{0.9em}} >{\raggedright\arraybackslash}X @{\hspace{0.9em}} >{\raggedleft\arraybackslash\bfseries}p{2.4em}@{}}
    #1 & #2 & #3
    \end{tabularx}
    \end{minipage}%
  }%
  \par\addvspace{0.32em}%
}

\title{\textbf{In the middle of the bridge}}
\author{\normalsize W24 (2024) --- English translation}
\date{}
\begin{document}
\maketitle

In this problem, we study electrical circuits containing a nonlinear element, the current-voltage characteristic of which is shown in Fig. 1. The current-voltage characteristic of a nonlinear element is described by two sections: $I=I_0\cdot(U/U_0)$ at $U\leq U_0$, where $I_0=10~\text{mA}$ and $U_0=5~\text{V}$. $I=I_0=const$ at $U\geq U_0$.

\begin{figure}[H]
    \centering
    \includegraphics[width=0.72\textwidth]{T2_EN_assets/fig1.jpg}
    \caption{Figure 1.}
\end{figure}

\section*{Part A. DC Circuit (3.0 points)}

Figure 2 shows a diagram of an electrical circuit consisting of resistors with resistances $R$, $2R$, $4R$, $6R$ and $8R$, an ideal constant voltage source $\mathcal{E}$ and the described nonlinear element. This part of the problem is devoted to studying the dependence of the current through a nonlinear element depending on the source voltage.

\begin{figure}[H]
    \centering
    \includegraphics[width=0.72\textwidth]{T2_EN_assets/fig2.jpg}
    \caption{Figure 2.}
\end{figure}

\problembox{A1}{Since without a section of the circuit between nodes $A$ and $B$ the electrical circuit consists entirely of constant voltage sources and resistors, the dependence of the current $I_{AB}$ flowing from node $A$ to node $B$ on the voltage $U_{AB}=\varphi_A-\varphi_B$ in this section of the circuit can be represented in a form that does not depend on the element connected between nodes $A$ and $B$: $$I_{AB}=\textbackslash{}cfrac{\textbackslash{}mathcal{E}_0-U_{AB}}{r_0}{.} $$ $\mathcal{E}_0$ and $r_0$.    $\mathcal{E}$ and $R$.}{1.50}

\problembox{A2}{Obtain the dependence of the current strength $I_{AB}$ on the value $\mathcal{E}$. Express your answer using $\mathcal{E}$, $R$, $U_0$ and $I_0$.}{0.50}

\problembox{A3}{Obtain the dependence of the power $P_\text{n.e}$ released on the nonlinear element on the value $\mathcal{E}$. Express your answer using $\mathcal{E}$, $R$, $U_0$ and $I_0$.}{0.50}

\problembox{A4}{Construct a high-quality graph of the dependence $P_\text{n.e}(\mathcal{E})$. Indicate all the characteristic points on the graph and indicate their corresponding values.}{0.50}

\section*{Part B. Electric Circuit Containing a Capacitor (5.0 points)}

The figure shows an electrical circuit consisting of an ideal voltage source $\mathcal{E}$, a resistor $R = 500~Ohm$, a capacitor with a capacitance $C = 100~\mu F$ and a nonlinear element connected in series, as well as a switch $K$. Initially, the key is open and the capacitor is not charged. At the initial time $t = 0$ the key is closed.

\begin{figure}[H]
    \centering
    \includegraphics[width=0.72\textwidth]{T2_EN_assets/fig3.jpg}
    \caption{Figure 3.}
\end{figure}

Let the emf of the source be equal to $\mathcal{E} = 6~\text{V}$.

\problembox{B1}{Obtain the dependence of the current $I_C(t)$ through the capacitor on time $t$. Express your answer in terms of $\mathcal{E}$, $R$, $U_0$, $I_0$, $C$ and $t$.}{0.80}

\problembox{B2}{What amount of heat $Q_\text{n.e}$ will be released on the nonlinear element during the establishment of constant current? Express your answer in terms of $\mathcal{E}$, $R$, $U_0$, $I_0$ and $C$. Find the numerical value.}{0.60}

\problembox{B3}{After what time $\tau$ the amount of heat released on the nonlinear element will reach $80\%$ from the value $Q_\text{n.e}$? Express your answer in terms of $R$, $U_0$, $I_0$, and $C$. Find the numerical value.}{0.60}

Now the emf of the source is $\mathcal{E} =12~\text{V}$.

\problembox{B4}{Obtain the dependence of the current $I_C(t)$ through the capacitor on time $t$. Express your answer in terms of $\mathcal{E}$, $R$, $U_0$, $I_0$, $C$ and $t$.}{1.50}

\problembox{B5}{What amount of heat $Q_\text{n.e}$ will be released on the nonlinear element during the establishment of constant current?  Find the numerical value with an accuracy of at least $1\%$.}{1.50}

\vspace{1em}
\noindent\textit{Source:} \url{https://pho.rs/p/3553}

\end{document}